398
\qquad\qquad\qquad\qquad\qquad\qquad\qquad\qquad\qquad\qquad 6 \text{ Integration}
\section{6.5.2 Convergence of an Improper Integral}
\subsection{a. The Cauchy Criterion}
\text{By Definition 3, the convergence of the improper integral (\ref{eq:6.71}) is equivalent to the}
\text{existence of a limit for the function}
\begin{equation}
F(b) = \int_{a}^{b} f(x) dx
\label{eq:6.74}
\end{equation}
\text{as } b \to \omega, b \in [\alpha, \omega[.
\text{This relation is the reason why the following proposition holds.}
\begin{proposition} 2 \text{ (Cauchy criterion for convergence of an improper integral) If the}
\text{function } x \mapsto f(x) \text{ is defined on the interval } [a, \omega[ \text{ and integrable on every closed}
\text{interval } [a, b] \subset [a, \omega[, \text{ then the integral } \int_{a}^{\omega} f(x)dx \text{ converges if and only if for}
\text{every } \varepsilon > 0 \text{ there exists } B \in [a, \omega[ \text{ such that the relation}
\begin{equation*}
\left| \int_{b_1}^{b_2} f(x)dx \right| < \varepsilon
\end{equation*}
\text{holds for any } b_1, b_2 \in [a, \omega[ \text{ satisfying } B < b_1 \text{ and } B < b_2.
\end{document}